\documentclass{article}

\usepackage{float}
\usepackage{xcolor}
\usepackage{parskip}
\usepackage{eulervm}
\usepackage{graphicx}
\usepackage{setspace}
\usepackage{geometry}
\usepackage{mathpazo}
\usepackage{microtype}
\usepackage[T1]{fontenc}
\usepackage[utf8]{inputenc}
\usepackage{amsfonts, amsmath, amssymb}
\usepackage[affil-it, auth-lg]{authblk}
\usepackage[colorlinks=true, allcolors=blue]{hyperref}
\usepackage[backend=biber, style=chicago-authordate, autocite=inline]{biblatex}
\usepackage[nameinlink, capitalize]{cleveref}   % Cleveref package must be loaded after all package

\geometry{paper=a4paper, portrait=true, layout=a4paper}

\title{The flaping of elephent's ear: where nature turns into engineering}

\author[1]{Mahriz Hossain}
\author[2]{Abdullah al Azmi}
\author[3]{Shadman Shamin Prionto}

\affil[1-3]{Department of Mechanical and Production Engineering, Islamic University of Technology (IUT), Bangladesh}

\addbibresource{references/references.bib}

\begin{document}

\maketitle

\begin{abstract}

	Nature has always surpised us. Human explored and it's beauty never end. Elephent, a mamalia kingdom's animal which is famous for it's large body. Though their large body, it's been a curosity of human being why there ear is so big while large animal like chil or wheel doesn't have. From the curosity, human start explored and upto current studies, it is evident that they use the ear for both purpose: listening and heat or thermoregulation. In this review, we expolored how the thermoregulation of elephent ear happened and what are the governing mechanism that drive this heat dissipitaion. Lastly, we mentioned what are the limitations current technologies and study have and discussed of the rooms for future improvements.

\end{abstract}

\section{Introduction}

A scientific group performed a study on african elephent where they found that the ear blood flow can constitue a substantial and controllable avenue of heat loss \autocite{dominguez-olivaAnatomicalPhysiologicalBehavioral2022}.

Further heat loss was most likely assisted by ear flapping which (although not quantified) indeed appeared to be much more frequent in elephants at the Thai study site \autocite{weissenbockTakingHeatThermoregulation2012}.

Likewise, the nonevaporative strategies mentioned by Dunkin et al (2013), such as the movement of the ears or shade seeking, are other tools that favor heat loss \autocite{dominguez-olivaAnatomicalPhysiologicalBehavioral2022}.

Ears are highly vascularized structures with movements that promote heat loss through convection (Mikota 2006).

Ear shape also affects thermoregulation as African species have up to twice the ear surface area (Dunkin et al 2013).

\section{Governing mechanisms}



\section{Discussion}



\end{document}
